%% ma: L11-B16-0034
%% lop: 11
%% chuong: 5
%% bai: 16
%% dang: 11.16.3
%% loai: TN4
%% muc: VD
%% dap_an: "B"
%% nguon: "tai-lieu-hoc-tap-mon-toan-lop-11-theo-chuong-trinh-giao-duc-pho-thong-2018.pdf (Lớp Toán Thầy Duy, Phạm Lê Duy, THPT Chu Văn An), Chương 3 Bài 2, Phần 1, Câu 22 (tr.366)"
%% kien_thuc: "Lớp 11 Bài 16 (giới hạn của hàm số)"
%% trang_thai: chinh-thuc
%% ghi_chu: "Đã duyệt 03/11/2026."
\begin{ex}
Cho hàm số $f(x)=\dfrac{(4x+1)^3(2x+1)^4}{(3+2x)^7}$. Tính $\lim\limits_{x\to-\infty}f(x)$.
\choice{$2$}{\True $8$}{$4$}{$0$}
\loigiai{\sloppy Chia tử và mẫu cho $x^7$: \[\dfrac{\left(4+\frac1x\right)^3\left(2+\frac1x\right)^4}{\left(\frac3x+2\right)^7}\to\dfrac{4^3\cdot2^4}{2^7}=8.\]}
\end{ex}
